The phenomenon is as follow:
% FIGURE NEEDED: Earth and Moon on a line inclined at angle theta to the
% ecliptic plane, with the centre of mass CM on the ecliptic and the Earth's
% perpendicular displacement x marked (2012 theory solutions, item 12, p. S5).

The Center of Mass of the Earth--Moon system (CM) lies on the plane of the
ecliptic, so as the Lagrangian points of its orbit around the Sun. Let be $D$
the distance between Earth and Moon and $x$ the distance between Earth and
the CM. We then have:
\[
M_\oplus \cdot x = M_{\mathrm{Moon}} \cdot (D - x)
\ \therefore\
x = \left(\frac{M_{\mathrm{Moon}}}{M_\oplus + M_{\mathrm{Moon}}}\right)D
\]
\hfill(2 pt)

The Earth's displacement perpendicular to the ecliptic is $A = 2x\sin\theta$,
where $\theta = 5^\circ$ is the inclination of the plane of Moon's orbit
relative to the plane of ecliptic. Calculating:
\[
A = \frac{2D\,M_{\mathrm{Moon}}}{M_\oplus + M_{\mathrm{Moon}}}\sin 5^\circ
 = 810\ \mathrm{km}
\]
\hfill(4 pt)

The telescope resolution is given by
\[
r'' = 1.22\frac{\lambda}{\delta} = \frac{A}{d}
\]
\hfill(3 pt)

where $\delta$ is the diameter of the objective lens and $d$ is the distance
to the telescope to Earth. But, at L4(L5), Earth and Sun forms an equilateral
triangle, so $d = 1$\,UA. Using $\lambda = 400$\,nm (that is, the visible
wavelength that allows the smallest telescope), we find
$\delta = 9.01\ \mathrm{cm}$.
\hfill(1 pt)
